
\documentclass{article}
\usepackage[utf8]{inputenc}
\usepackage[russian]{babel}
\usepackage{geometry}
\usepackage{hyperref}
\usepackage{booktabs}
\usepackage{longtable}

\title{Филологический отчет RuWritingStyles}
\author{Агентская Лаборатория v2.1}
\date{\today}

\begin{document}

\maketitle

\section{Метаданные}
\begin{itemize}
    \item \textbf{ID запуска:} phase1-council-example
    \item \textbf{Профиль:} None
    \item \textbf{Модель:} None
    \item \textbf{Длительность:} 0.00 сек.
\end{itemize}

\section{Методологический Компас}
Результаты анализа принадлежности к научным школам:
\begin{itemize}
    \item \textbf{{Leningrad}}: 0.80
    \item \textbf{{Moscow}}: 0.20
\end{itemize}

\section{Когнитивная разметка (Bloom)}
Распределение когнитивной нагрузки в решениях Совета:
\begin{itemize}
    \item \textbf{{Analyze}}: 4
    \item \textbf{{Evaluate}}: 4
\end{itemize}

\section{Результаты аудита (Совет)}
\begin{longtable}{@{}p{0.2\textwidth}p{0.7\textwidth}@{}}
\toprule
\textbf{ID находки} & \textbf{Решение и Обоснование} \\
\midrule
finding-001 & Mock council keeps placeholder findings informational. \\ \hline
finding-001 & Mock council keeps placeholder findings informational. \\ \hline
finding-001 & Mock council keeps placeholder findings informational. \\ \hline
finding-001 & Mock council keeps placeholder findings informational. \\ \hline

\bottomrule
\end{longtable}

\section{Аудит методологической предвзятости}
Результаты анализа беспристрастности Совета:
\begin{itemize}
    \item \textbf{Оценка смещения:} 1/10
    \item \textbf{Тип смещения:} none
    \item \textbf{Критика:} Mock critique.
\end{itemize}

\section{Библиографический конкорданс}
Ссылки на академические коллекции:
В данном запуске использованы ссылки на коллекции Зализняка и Тронского.

\section{Научная обоснованность (Citations)}
Результаты проверки цитирований:
\begin{itemize}
    \item \textbf{{Всего цитат:}} 4
    \item \textbf{{Верифицировано:}} 4
    \item \textbf{{Не в библиографии:}} 0
\end{itemize}

\section{Литература}
Список литературы по ГОСТ Р 7.0.100--2018 (кириллица перед латиницей):
\begin{enumerate}
    \item Зализняк А. А. Древненовгородский диалект. — 2-е изд., перераб. с учетом материала находок 1995–2003 гг. — М. : Языки славянской культуры, 2004. — 872 с.
    \item Казанский Н. Н. Проблемы комментирования и сопоставления в русско-латинском параллельном корпусе // Индоевропейское языкознание и классическая филология. — 2025. — Т. 29. — С. 857–873.
    \item Парибок А. В. Уровень смысла в текстах шастр: концепты или понятия? // Шабдапракаша. Зографский сборник I. — СПб., 2011. — С. 80–86.
    \item Эрман В. Г. Махабхарата. Книга шестая. Бхишмапарва / пер. с санскрита, предисл. и коммент. В. Г. Эрмана. — М. : Ладомир, 2009.
\end{enumerate}

\end{document}
